\section*{Conclusion}

We have introduced ArchaicPainter (\AP{}), a Hidden Markov Model that adapts the
tract-length transition structure of the Li \& Stephens framework to archaic hominin
introgression detection, with hidden states representing ancestry categories rather
than individual reference haplotypes.
By encoding coalescent match-length statistics in the emission and transition
probabilities of a three-state HMM, \AP{} exploits an $\sim$11-fold
coalescent tract-length contrast between introgressed and non-introgressed segments
that all existing density-based approaches leave unexploited.
Across 50 simulation replicates with known ground truth, this architectural
choice achieves F1 $= 0.480 \pm 0.095$ (AUPRC $= 0.577$) across 50 simulation
replicates, validating the discriminative value of encoding match-length structure.

On real 1000~Genomes Project chromosome~21 data, \AP{} identifies
$\sim$1.1\% (CEU) and $\sim$1.0\% (CHB) per-haplotype Neanderthal-like
haplotype coverage at a median segment resolution of $\sim$28--57~kb,
broadly consistent in order of magnitude with published genome-wide estimates of 1--2\%.
Because real-data calls use a heuristic positive-only log-score (not a calibrated
probability model), these fractions are indicative estimates rather than
posterior-derived quantities; they also reflect lower recall for short segments
under sparse informative-site coverage ($\sim$27 SNPs/Mb on chr21).
Functional annotation of detected regions recovers known adaptive introgression
targets including the complete type-I and type-II interferon receptor gene cluster
in an unsupervised manner, corroborating published evidence for adaptive retention
of Neanderthal antiviral immunity haplotypes~\citep{dannemann2017contribution,
enard2018human} and illustrating the method's ability
to generate biologically interpretable output for downstream GWAS integration.

The results motivate two conclusions for the field.
First, haplotype match-length is a strong and previously underutilized signal
for archaic introgression: methods designed to capture it --- whether through
the LS framework or related approaches --- warrant serious development alongside
existing density-based tools.
Second, per-haplotype resolution, as opposed to per-individual diploid resolution,
opens qualitatively new analyses: phased introgression catalogs can be directly
connected to phased GWAS loci, enabling haplotype-level hypothesis testing about
archaic allele contributions to complex traits that diploid-average methods cannot support.

Limitations of the present work --- chromosome~21 analysis only, no direct
benchmark against HMMix or DAIseg, proof-of-concept attribution for Denisovan
ancestry --- are explicitly documented and define a concrete roadmap for
follow-on work: genome-wide analysis, head-to-head benchmarking, integration of
real Denisovan reference data, and SparsePainter-based scaling to biobank cohorts.
\AP{} is released as open-source software at \url{https://github.com/EvoClaw/ArchaicPainter}
with fully reproducible analysis pipelines.
We anticipate that per-haplotype resolution archaic introgression maps,
integrated with biobank-scale GWAS, will reveal new dimensions of how Neanderthal
and Denisovan haplotypes shape human phenotypic diversity and disease risk today.
